\SourceLeaf{187}und jetzt beinahe überzeugt sei, dass ohne irgend ein datum der Euklidische Satz nicht zu beweisen sei, dass es ihm auch nicht unwahrscheinlich sei, dass unsere Geometrie nur ein Kapitel einer allgemeinern sei. Ich erzählte ihm darauf, wie Sie vor einigen Jahren öffentlich geäussert hätten, dass man seit Euklids Zeiten im Grunde hiermit nicht weiter gekommen sei; ja dass Sie gegen mich mehrmals geäussert hätten, wie Sie durch vielfältige Beschäftigung mit diesem Gegenstand, auch nicht zum Beweise von der Absurdität einer solchen Annahme gekommen seien. --- Als er mir darauf das verlangte Buch schickte, lag der beigehende Zettel bei, und er bat mich kurz darauf (Ende December) mündlich, Ihnen doch gelegentlich diesen seinen Zettel beizuschliessen, und Sie in seinem Namen zu ersuchen, gelegentlich ihm Ihr Urtheil über seine Ideen wissen zu lassen.

Das Buch selbst hat, abgesehen von allem übrigen, das angenehme, dass eine reichhaltige Literatur des Gegenstandes sich darin findet; welche er auch, wie er mir sagt, ferner zu sammeln nicht abgelassen hat. .\ .\ .\ .\ .\}

\bigskip
\noindent\hrulefill
\bigskip

\Needspace{10\baselineskip}
\begin{center}[\emph{\textls{Beilage:} Notiz von} \textsc{Schweikart}. \emph{Marburg, December 1818.}]\end{center}

\{Es gibt eine zwiefache Geometrie, --- eine Geometrie im engern Sinn --- die \textls{Euklidische}; und eine astralische Grössenlehre.

Die Dreiecke der letztern haben das Eigene, dass die Summe der drei Winkel nicht zwei Rechten gleich ist.

\textls{Diess} vorausgesetzt, lässt es sich auf das strengste beweisen:

a) dass die Summe der 3 Winkel in dem Dreieck \textls{kleiner} als 2 Rechte sei;

b) dass diese Summe immer kleiner werde, je mehr Inhalt das Dreieck umfasst;

c) dass die Höhe eines gleichschenkligen rechtwinkligen Dreiecks zwar immer zunimmt, je mehr man die Schenkel verlängert, dass sie aber eine gewisse Linie, die. ich die \textls{Constante} nenne, nicht übersteigen könne.

